\section{Classical local algorithms}\label{sec:classical}
This section gives classical algorithms that estimate $q(H,b)$ in the
access model of Section~\ref{sec:models}. The ingredients are standard:
Chebyshev residual polynomials, local-walk evaluation of sparse matrix
polynomials, and importance sampling. Gharibian--Le Gall
\citep{GharibianLeGall2022} already estimate sparse-polynomial overlaps;
Montanaro--Shao \citep{MontanaroShao2023} and
Cifuentes et al.\ \citep{CifuentesEtAl2024} study matrix-function entry
estimation. Our version (Theorem~\ref{thm:classical-scalar}) guarantees
relative error for the positive inverse quadratic form and uses positivity
to reduce the statistical factor from $\kappa^2$ to $\kappa$. We then give
two variants for comparison with the optimization constructions: a
solution sampler (a stronger output) and an estimator for nonsymmetric
matrices.

\subsection{A relative inverse polynomial}
\begin{lemma}[Chebyshev residual]\label{lem:residual}
For $K>1$ and $0<\eta<1$, there is a real polynomial $P$ of degree at most
$m-1$, where
\begin{equation}\label{eq:degree}
 m=\left\lceil\frac{\log(2/\eta)}{
       \log((\sqrt K+1)/(\sqrt K-1))}\right\rceil
       =O(\sqrt K\log(2/\eta)),
\end{equation}
such that $|1-\lambda P(\lambda)|\leq\eta$ for
$\lambda\in[K^{-1},1]$. Consequently, for
$K^{-1}I\preceq H\preceq I$,
\begin{equation}\label{eq:relative-polynomial}
 (1-\eta)H^{-1}\preceq P(H)\preceq(1+\eta)H^{-1},\qquad
 \norm{I-HP(H)}\leq\eta.
\end{equation}
For $K=1$, one may take $P=1$.
\end{lemma}
\begin{proof}
Put $a=K^{-1}$, $c=(1+a)/(1-a)$, and
$x(\lambda)=(1+a-2\lambda)/(1-a)$. If $T_m$ is the degree-$m$
Chebyshev polynomial, set
\[
 r_m(\lambda)=\frac{T_m(x(\lambda))}{T_m(c)},\qquad
 P(\lambda)=\frac{1-r_m(\lambda)}{\lambda}.
\]
Since $x(0)=c$, the numerator vanishes at zero, so the quotient is a
polynomial. On $[a,1]$, $|T_m(x(\lambda))|\leq1$. With
$\theta=(\sqrt K-1)/(\sqrt K+1)$, the identity
$T_m(c)=(\theta^{-m}+\theta^m)/2$ gives
$|r_m(\lambda)|\leq2\theta^m$. The displayed choice of $m$ makes this at
most $\eta$. Finally
$\log((1+u)/(1-u))\geq2u$ for $u=1/\sqrt K$ proves the degree estimate;
the matrix inequalities follow by diagonalizing $H$.
\end{proof}

\begin{lemma}[Local evaluation]\label{lem:local}
If $H$ is $d$-row-sparse, one coordinate of $P(H)v$, for a specified
degree-$k$ polynomial $P$, can be evaluated with
$(d+1)^{O(k+1)}$ sparse and coordinate queries and arithmetic operations.
If row and column sparsity are both bounded by $d$, a full row or column of
$P(H)$ can be materialized in the same bound. More generally, a polynomial
$q(A^*A)A^*$ of degree $k$ in $A^*A$ has this property under row and column
access to $A$, with walk length at most $2k+1$.
\end{lemma}
\begin{proof}
Expand a requested coordinate of $H^jv$ into matrix walks of length $j$
starting from that coordinate. Each step has at most $d$ choices. Summing
for $0\leq j\leq k$ costs at most
$O((k+1)^2(d+1)^{k+1})$ operations, including multiplying path weights and
aggregating equal endpoints. The same expansion without the terminal
coordinate of $v$ materializes the corresponding row; reverse walks
materialize a column. One can avoid sorting by retaining the walk list and
aggregating equal endpoints by comparisons in quadratic time, still within
the stated bound. For $q(A^*A)A^*$ the walks alternate between column and
row queries to $A$. No length-$N$ vector is formed.
\end{proof}

\subsection{Positive quadratic forms}
\begin{lemma}[Positive-form second moment and the Kantorovich inequality]\label{lem:kantorovich}
If $P\succ0$ has spectrum in $[\alpha,\beta]$ and $b\ne0$, sample
$I\sim p_b$ and define
\begin{equation}\label{eq:estimator}
                    X=\norm b^2\frac{(Pb)_I}{b_I}.
\end{equation}
Then
\begin{equation}\label{eq:second-moment}
 \E X=b^*Pb,\qquad
 \frac{\E|X|^2}{(b^*Pb)^2}
 \leq\frac{\norm b^2\norm{Pb}^2}{(b^*Pb)^2}
 \leq\frac{(\alpha+\beta)^2}{4\alpha\beta}.
\end{equation}
The last inequality is the classical Kantorovich inequality in an equivalent
form; see \citet[Eq.~(1.1)]{Lin2013}. It is sharp.
\end{lemma}
\begin{proof}
The sampling identity follows from $|b_i|^2/b_i=\overline b_i$ on the
support of $b$. The second moment is
$\norm b^2\sum_{i:b_i\ne0}|(Pb)_i|^2$, which is bounded by the expression
shown. Under the spectral weights of $b/\norm b$, let $Y$ be an eigenvalue
of $P$ and $t=\E Y$. Since
$(Y-\alpha)(Y-\beta)\leq0$,
\[
 \frac{\E Y^2}{(\E Y)^2}
 \leq\frac{(\alpha+\beta)t-\alpha\beta}{t^2}
 \leq\frac{(\alpha+\beta)^2}{4\alpha\beta}.
\]
The last maximum occurs at $t=2\alpha\beta/(\alpha+\beta)$.
Equality is attained by $P=\diag(\alpha,\beta)$ with squared coordinates
of normalized $b$ equal to $\beta/(\alpha+\beta)$ and
$\alpha/(\alpha+\beta)$. This proves sharpness for this estimator, without
asserting an algorithm-independent lower bound.
\end{proof}

\begin{theorem}[Relative scalar estimation]\label{thm:classical-scalar}
Suppose $H$ satisfies \eqref{eq:spectral-promise} and is $d$-row-sparse,
and that row-sparse access to $H$ and $\SQ(b)$ are available. There is a randomized
algorithm returning a relative-$\epsilon$ estimate of $q(H,b)$ with failure
probability at most $\zeta$, using
\begin{equation}\label{eq:classical-bound}
 O\!\left(\kappa\epsilon^{-2}\log(1/\zeta)
       (d+1)^{O(\sqrt\kappa\log(2/\epsilon))}\right)
\end{equation}
queries and ideal arithmetic operations. Full matrix $\SQ$ is not needed.
\end{theorem}
\begin{proof}
For $\kappa=1$, $H=I$ and the norm oracle returns $q=\norm b^2$ exactly.
Otherwise choose the polynomial in Lemma~\ref{lem:residual} with
$\eta=\epsilon/4$ and put $Q=b^*P(H)b$. Then
$|Q-q|\leq\eta q$. The eigenvalues of $P(H)$ lie in
$[1-\eta,\kappa(1+\eta)]$, so their ratio is at most
$9\kappa/7$, because $\eta\leq1/8$. Lemma~\ref{lem:kantorovich} gives
$\operatorname{Var}(\operatorname{Re}X)\leq C\kappa Q^2$ for an absolute
constant $C$.

Divide independent samples into an odd number
$g=O(\log(1/\zeta))$ of groups, each of size
$h=\lceil64C\kappa/\epsilon^2\rceil$. Chebyshev's inequality bounds by
$1/4$ the probability that a group mean differs from $Q$ by more than
$\epsilon Q/4$. The median of the group means fails only if at least half
the groups fail; the binomial Chernoff bound makes this probability at most
$\zeta$ for a suitable absolute constant in $g$. On success,
\[
 |\widehat q-q|
 \leq(\epsilon/4)Q+(\epsilon/4)q
 \leq(\epsilon/4)(2+\epsilon/4)q<\epsilon q.
\]
Each sample requires one local polynomial evaluation, covered by
Lemma~\ref{lem:local}; multiplying its cost by $gh$ proves the result.
\end{proof}

Standard inverse-form quadrature uses full matrix--vector products
\citep{LiSraJegelka2016}; the algorithm above explores only neighborhoods
of sampled coordinates. The price is cost exponential in the degree: the
local factor is $\exp(O(\sqrt\kappa\log(d+1)\log(2/\epsilon)))$, so the
bound is not polynomial in $\kappa$.

\begin{corollary}[Bilinear forms]\label{cor:bilinear}
With the same matrix assumptions, $\SQ(a)$ and coordinate access to $b$
suffice to estimate $a^*H^{-1}b$ to additive error
\[
       \epsilon\sqrt{(a^*H^{-1}a)(b^*H^{-1}b)}
\]
with the complexity in \eqref{eq:classical-bound}. If the public promise
$|a^*H^{-1}b|\geq\gamma\sqrt{(a^*H^{-1}a)(b^*H^{-1}b)}$ holds for
$0<\gamma\leq1$, replacing $\epsilon$ by $\gamma\epsilon$ gives relative
error $\epsilon$.
\end{corollary}
\begin{proof}
Write $A=a^*H^{-1}a$ and $B=b^*H^{-1}b$. The zero-vector cases return zero.
The residual representation gives
\[
 |a^*(P(H)-H^{-1})b|\leq\eta\sqrt{AB}.
\]
Indeed, $P(H)-H^{-1}=H^{-1/2}E H^{-1/2}$ for a Hermitian $E$ with
$\norm E\leq\eta$. Sampling from $p_a$, the random variable
$X=\norm a^2(P(H)b)_I/a_I$ has mean $a^*P(H)b$ and second moment at most
$\norm a^2\norm{P(H)b}^2$. The pointwise spectral inequality
$P(H)^2\preceq\kappa(1+\eta)^2H^{-1}$ and $\norm a^2\leq A$ bound this
by $\kappa(1+\eta)^2AB$. Apply the group-mean argument separately to real
and imaginary parts, with tolerances $\epsilon\sqrt{AB}/4$ and failure
probabilities $\zeta/2$. The sample count depends on the known upper bound
for the ratio of variance to tolerance squared, so it does not require
knowing $A$ or $B$. Taking $\eta=\epsilon/4$ proves the claim.
\end{proof}

\subsection{Sampling a sparse inverse solution}\label{subsec:sampling}
Scalar estimation does not require a solution sampler. For comparison,
the following elementary rejection scheme gives this stronger output.
\begin{lemma}[Sparse-transform rejection sampler]\label{lem:rejection}
Let $P$ be invertible, with at most $R$ nonzeros per row and column and
singular values in $[a_P,b_P]$. Given $\SQ(v)$ and explicit
materialization of the sparse rows and columns of $P$, one can sample exactly from $p_{Pv}$ using at most
$R(b_P/a_P)^2$ trials in expectation. Each trial materializes at most one
row and one column and uses $O(R)$ further coordinate operations.
\end{lemma}
\begin{proof}
Let $c_j^2=\sum_i|P_{ij}|^2$ and
$S_i=\sum_j|P_{ij}v_j|^2$. Sample $j\sim p_v$, accept with probability
$c_j^2/b_P^2$, and, if accepted, sample $i$ from column $j$ with probability
$|P_{ij}|^2/c_j^2$. Compute $y_i=(Pv)_i$ and $S_i$ from row $i$, and accept
$i$ with probability $|y_i|^2/(RS_i)$. The event $S_i=0$ has zero
probability at this stage and can be defined to reject. Column norm bounds
and Cauchy--Schwarz show both acceptance probabilities lie in $[0,1]$.
Summing over $j$, the probability of outputting $i$ in one raw trial is
\[
 \sum_j\frac{|v_j|^2}{\norm v^2}\frac{c_j^2}{b_P^2}
         \frac{|P_{ij}|^2}{c_j^2}\frac{|y_i|^2}{RS_i}
       =\frac{|y_i|^2}{Rb_P^2\norm v^2}.
\]
Conditioning on acceptance gives $p_{Pv}$. Its total probability is at least
$a_P^2/(Rb_P^2)$, proving the expected trial count.
\end{proof}

\begin{theorem}[Solution sampling]\label{thm:solution-sampling}
Let $A\in\C^{N\times N}$ be invertible with
$\norm A\leq1$, $\sigma_{\min}(A)\geq K^{-1}$, and at most $d$ nonzeros
per row and column. Given row- and column-sparse access to $A$ and
$\SQ(b)$, an algorithm can provide coordinate access to a fixed vector
$\widetilde x$ satisfying
\begin{equation}\label{eq:solution-error}
 \norm{\widetilde x-A^{-1}b}\leq\epsilon\norm{A^{-1}b},
\end{equation}
together with exact samples from $p_{\widetilde x}$, at expected cost
\begin{equation}\label{eq:sampling-cost}
 K^2(d+1)^{O(K\log(2/\epsilon))}
\end{equation}
per sample. For positive definite $A$, this improves to
$K^2(d+1)^{O(\sqrt K\log(2/\epsilon))}$.
\end{theorem}
\begin{proof}
Apply Lemma~\ref{lem:residual} to $B=A^*A$, whose spectrum is contained
in $[K^{-2},1]$, and let $q$ be its inverse polynomial with residual
$r=1-\lambda q(\lambda)$ bounded by $\epsilon$. Put
$P=q(A^*A)A^*$ and $\widetilde x=Pb$. If $x=A^{-1}b$, then
\[
                  Pb-x=-r(A^*A)x,
\]
which proves \eqref{eq:solution-error}. The singular value decomposition of
$A$ shows that the singular values of $P$ are
$|1-r(\sigma^2)|/\sigma$; thus they lie in
$[1-\epsilon,K(1+\epsilon)]$ and $\kappa(P)\leq3K$.
Lemma~\ref{lem:local} materializes each needed row or column in cost
$W=(d+1)^{O(K\log(2/\epsilon))}$, and also bounds its nonzero count by
$W$. Lemma~\ref{lem:rejection} therefore costs $O(K^2W^2)$ per sample,
which has the stated form. For positive definite $A$, use $P=q(A)$ directly
with degree $O(\sqrt K\log(2/\epsilon))$. When $K=1$, the general matrix
is unitary and $P=A^*$ is exact; the same rejection proof applies.
\end{proof}

For fixed sparsity and accuracy, this upper bound implies that an $N^c$
classical lower bound with the same access and output requirement needs $K=\Omega(\log N)$ in general and
$K=\Omega(\log^2N)$ in the positive definite case. These are necessary
conditions; we do not claim hardness for every matrix at those scales.
Andoni--Krauthgamer--Pogrow explicitly ask how to construct an
$\ell_2$ sampler for a solution from an $\ell_2$ sampler for its right-hand
side \citep[Section~1.4]{AndoniKrauthgamerPogrow2019}; their main results
instead give coordinate estimates under different hypotheses and error
guarantees. Theorem~\ref{thm:solution-sampling} states the sampling guarantee and
rejection cost explicitly.

\subsection{Nonsymmetric inverse forms and optimization values}
\begin{theorem}[Norm-sensitive inverse overlap]\label{thm:general-overlap}
Let $M$ satisfy the general hypotheses of Theorem~\ref{thm:solution-sampling}.
Given $\SQ(a)$ and coordinate access to $e$, the scalar $a^*M^{-1}e$ can be
estimated to additive error $\epsilon\norm a\norm e$, with failure
probability $\zeta$, at cost
\begin{equation}\label{eq:general-overlap}
 O\!\left(R_e^2\epsilon^{-2}\log(1/\zeta)
 (d+1)^{O(K\log(2R_e/\epsilon))}\right),
\end{equation}
where $1\leq R_e\leq K$ is any supplied upper bound on
$\norm{M^{-1}e}/\norm e$. Without a sharper bound, take $R_e=K$.
Here $0<\epsilon\leq R_e/2$, which extends the default range
$\epsilon\leq1/2$ of Section~\ref{sec:models} when $R_e>1$.
The modulus $|a^*M^{-1}e|$ admits the same additive accuracy.
\end{theorem}
\begin{proof}
Take the normal-equation residual polynomial from the preceding proof with
$\eta=\epsilon/(4R_e)\leq1/8$, and put $P=q(M^*M)M^*$.
Its error satisfies
\[
 \norm{Pe-M^{-1}e}\leq\eta R_e\norm e,\qquad
 \norm{Pe}\leq(1+\eta)R_e\norm e.
\]
The estimator $X=\norm a^2(Pe)_I/a_I$, with $I\sim p_a$, has mean
$a^*Pe$ and second moment at most
$(1+\eta)^2R_e^2\norm a^2\norm e^2$. The polynomial bias is at most
$\epsilon\norm a\norm e/4$. Separate real and imaginary median-of-means
estimates give the claimed statistical error using
$O(R_e^2\epsilon^{-2}\log(1/\zeta))$ samples. Local evaluation costs
$(d+1)^{O(K\log(2R_e/\epsilon))}$. The inequality
$||z|-|\widehat z||\leq|z-\widehat z|$ proves the modulus claim.
\end{proof}
Corollary~\ref{cor:bilinear} normalizes the error by the
$H^{-1}$-norms of $a$ and $b$; Theorem~\ref{thm:general-overlap}
normalizes it by the Euclidean norms of $a$ and $e$. Neither gives a
relative estimate of an overlap that can be small through cancellation.
The parameter $R_e$ is useful when an optimization value equals
$|a^*M^{-1}e|$ and the construction makes the solution norm public.

\subsection{What finite precision changes}\label{subsec:precision}
\citet{LeGall2025} studies robust dequantization under approximate
length-squared sampling. The oracle bounds of this section assume exact
sampling and ideal arithmetic; they do not claim that approximate
sampling is free, and they give no bound on the bit precision an
implementation needs. The following proposition gives a
sufficient condition under which the estimator of
Theorem~\ref{thm:classical-scalar}, with approximate samples and inexact
arithmetic, keeps relative error $\epsilon$ at failure probability $2\zeta$.
\begin{proposition}[Robustness to approximate samples and arithmetic]\label{prop:precision}
Run the polynomial estimator of Theorem~\ref{thm:classical-scalar} with
$L=O(\kappa\epsilon^{-2}\log(1/\zeta))$ samples. Suppose the distribution
of each requested sample, conditional on all previous computation, is
within total-variation distance $\delta$ of $p_b$. Extend the estimator by the value zero on
indices with $b_i=0$, and have the approximate implementation return zero
on those indices. On the support of $b$, suppose every reported real
estimator value $\widehat X_i$ differs from its exact value
$\operatorname{Re}X_i$ by at most $\tau\norm b^2$.
There is a coupling with an exact-oracle execution under which the final
medians differ by at most $\tau\norm b^2$, except on an event of
probability at most $L\delta$. In particular, $\tau\leq\epsilon/4$ and
$\delta\leq\zeta/L$ retain relative error $\epsilon$ and failure
probability at most $2\zeta$ for the constants used above.
\end{proposition}
\begin{proof}
Successive maximal couplings make every sample agree with the exact-oracle
sample except with probability at most $L\delta$, by a conditional union
bound. An off-support approximate sample necessarily belongs to the
disagreement event. On agreement, every group mean changes by at most
$\tau\norm b^2$; the median is Lipschitz with constant one in the maximum
coordinate norm. Since $q\geq\norm b^2$, this is at most $\tau q$.
The proof of Theorem~\ref{thm:classical-scalar} leaves error at most
$17\epsilon q/32$, so the additional $\epsilon q/4$ fits within the
claimed tolerance.
\end{proof}

If the norm and the sampled coordinate $b_i\ne0$ are exact, the error
condition of Proposition~\ref{prop:precision} holds when the computed
polynomial coordinate satisfies
\[
 |\widehat{(P(H)b)_i}-(P(H)b)_i|\leq\tau|b_i|.
\]
Alternatively, interval arithmetic can bound the estimator error
directly, including its division and norm multiplication. For finite rational input data and a rational spectral bound, the
residual polynomial has rational coefficients. A finite local-walk sum
and its sampled nonzero denominator can therefore be evaluated to any
such tolerance, or exactly in rational arithmetic.
This shows implementability, not a dimension-independent bit bound;
such a bound must account for small sampled coordinates, input bit
lengths, coefficient growth, and the cost of approximate sampling.
In particular, boundedness of the polynomial on the spectral interval
does not by itself justify a uniform working precision for divisions by
arbitrarily small $b_i$. Proposition~\ref{prop:precision} is the
finite-precision claim used here; the other oracle bounds assume ideal
arithmetic.
